% TOP_PROBLEM: {"id": 3362, "title": "Birch & Swinnerton-Dyer conjecture", "category_id": 1, "category": {"id": 1, "name": "number_theory", "display_name": "Number Theory", "description": "Properties of integers, prime numbers, Diophantine equations.", "slug": "number-theory", "order_index": 1, "created_at": "2026-07-31T15:26:25.670Z"}, "status": "open", "problem_number": "OPG-37423", "set_id": 12}
% FIRST_POSTED: 2026-10-01
% ENTRY_KIND: statement_only
% UnsolvedMath ID 3362; source catalogue code OPG-37423
% !TeX program = lualatex
% Definition and sources only; this file is not an LLM solution attempt.
\documentclass[11pt,a4paper]{article}
\usepackage[margin=25mm]{geometry}
\usepackage{fontspec}
\setmainfont{lmroman10-regular.otf}[BoldFont=lmroman10-bold.otf,
 ItalicFont=lmroman10-italic.otf,BoldItalicFont=lmroman10-bolditalic.otf]
\usepackage{amsmath,amssymb,mathtools}
\usepackage{unicode-math}
\setmathfont{latinmodern-math.otf}
\AtBeginDocument{\let\setminus\smallsetminus}
\usepackage{xurl,hyperref}
\hypersetup{colorlinks=true,urlcolor=blue}
\setcounter{secnumdepth}{0}
\setlength{\parindent}{0pt}
\setlength{\parskip}{5pt}
\setlength{\emergencystretch}{3em}
\begin{document}
\section*{Birch \& Swinnerton-Dyer conjecture}\label{3362}
{\small UnsolvedMath ID 3362 · OPG-37423 · Number Theory · OpenGarden}
\subsection{Definitions and mathematical statement}
Let $K$ be a number field, a finite extension of $\mathbb Q$, and let
$E/K$ be an elliptic curve: a smooth projective genus-one curve with
a specified $K$-rational point. Its rational points form a finitely
generated abelian group $E(K)\cong\mathbb Z^r\oplus T$, with $T$
finite; $r$ is the Mordell--Weil rank.

For a finite place $v$ of good reduction, let $q_v$ be the size of
its residue field and put $a_v=q_v+1-\#E(\mathbb F_{q_v})$. The
Hasse--Weil $L$-function, initially defined in a right half-plane, is
\[
 L(E/K,s)=\prod_{v\text{ good}}
 (1-a_vq_v^{-s}+q_v^{1-2s})^{-1}
 \prod_{v\text{ bad}}(1-a_vq_v^{-s})^{-1}.
\]
At bad places, $a_v$ is $1$ for split multiplicative reduction, $-1$
for nonsplit multiplicative reduction, and $0$ for additive reduction.
The rank form of the conjecture asserts analytic continuation to a
neighborhood of $s=1$ and
\[
                    \operatorname{ord}_{s=1}L(E/K,s)=r.
\]
The order is the integer $m\ge0$ for which
$L(E/K,s)=(s-1)^m g(s)$ with $g$ holomorphic and $g(1)\ne0$.
The quantifier covers every elliptic curve over every number field.
This record concerns the rank equality; the stronger formula for the
leading Taylor coefficient is not included in its source statement.
\subsection{Short English statement}
For every elliptic curve over a number field, does its number of independent rational points equal the order of vanishing of its L-function at one?
\subsection{Sources}
\begin{itemize}
\item[S1] ULAM.AI and the UnsolvedMath Contributors, supplied problems.json, record 3362 (OPG-37423), OpenGarden collection. Catalogue identity and source transcription; CC BY 4.0.
\url{https://huggingface.co/datasets/ulamai/UnsolvedMath}.
\item[S2] Open Problem Garden, Birch \& Swinnerton-Dyer conjecture. Original problem and discussion; the mathematical formulation below is expanded from the supplied source record.
\url{http://www.openproblemgarden.org/op/birch_swinnerton_dyer_conjecture}.
\end{itemize}
\end{document}
